\documentclass[11pt]{article}

\usepackage[a4paper,margin=28mm]{geometry}
\usepackage{amsmath,amssymb,amsthm,mathtools}
\usepackage[hidelinks]{hyperref}
\usepackage{microtype}

\newtheorem{lemma}{Lemma}
\newcommand{\Z}{\mathbb Z}
\newcommand{\Rel}{\mathop{\rm Rel}}
\newcommand{\Sym}{\mathop{\rm Sym}}
\newcommand{\src}{\mathop{\rm src}}
\newcommand{\ev}{\mathop{\rm ev}}

\title{The Exact Step--7 Obstruction\\
  in the Notation of the Original Proof}
\author{}
\date{}

\begin{document}
\maketitle

\section{Purpose}

This note isolates the one mathematical assertion that is still missing from
the formal verification of Step~7.  It is not another version of the final
theorem, and it does not conceal Point~6 or any part of the reduction.  It is
the terminal-source realization of one external edge of the marked-row
calculation: the exceptional relation-on-the-left edge.

The assertion is most naturally a \emph{global statement about signed row
occurrences}.  A statement about a single homogeneous component of a single
relation is not sufficient and is generally false at the lower contexts.

\section{The manuscript setup}

Let
\[
       F=\bigoplus_{s=1}^{n+1}F_s
\]
be the truncated free metabelian Lie ring used in the proof, let
\(\ev:F\to L\) be its evaluation, and put
\[
       \Rel=\ker(\ev).
\]
Write
\[
 A_n=\bigoplus_{s=1}^{n}F_s,
 \qquad
 D_n=\operatorname{im}(\Rel\longrightarrow A_n),
 \qquad
 d_n:D_n\hookrightarrow A_n.
\]
The degree-one part of the terminal quadratic Koszul complex is
\[
 K_2(d_n)_1=D_n\otimes A_n,
\]
and its differential is
\[
 \partial_n(r\otimes x)=r\odot x\in\Sym^2 A_n.
\]

Fix the integral weighted PBW basis and the full-relation and generator
sections used for the terminal presentation.  If
\(z\in D_n\otimes A_n\), let
\[
       \src_n(z)\in F
\]
denote the complete one-factor PBW primitive of its literal
relation-on-the-left source placement.  This is the source primitive which
is later compared with the four Smith families.

Now choose the finite relative one-chain whose Chevalley--Eilenberg
differential is the governing element \(\Theta\in\Rel U(F)\), and run the
complete occurrence-labelled marked-row calculation of the proof.  At the
mark-one truncation wall, let \(\mathcal E\) be the signed set of occurrences
having all three properties:
\begin{enumerate}
\item the bracket context is empty;
\item the full relation has leading weight one;
\item the ordered spectator word has at least two factors.
\end{enumerate}
Thus an occurrence \(e\in\mathcal E\) has a coefficient
\(m_e\in\Z\), a \emph{full} triangular relation
\(\rho_e\in\Rel\), and an ordered word
\[
       X_e=x_{e,1}\cdots x_{e,r_e},
       \qquad r_e\geq2.
\]
No homogeneous piece of \(\rho_e\) is being declared a relation.

Let \(P_1\) and \(P_2\) be the complete factor-one and factor-two PBW
projections in \(U(F)\), and let \(\pi_n:F\to A_n\) be the prefix
projection.  Define the two reads of the exceptional relation-left edge by
\begin{align*}
 B_{\mathcal E}
   &=\sum_{e\in\mathcal E}m_e\,
       \Sym^2(\pi_n)P_2\bigl(\iota(\rho_e)X_e\bigr)
       \in\Sym^2A_n,\\
 P_{\mathcal E}
   &=\sum_{e\in\mathcal E}m_e\,
       P_1\bigl(\iota(\rho_e)X_e\bigr)
       \in F.
\end{align*}
The relation is on the left in both formulas.  These are exactly the upper
factor-two edge and its complete PBW primitive; the coefficients are the
coefficients of labelled occurrences, not coefficients obtained by expanding
a relation component.

\section{The missing assertion}

\begin{lemma}[Exceptional left-edge completion]\label{lem:missing}
For the signed exceptional occurrence family produced by the complete
marked-row trace of the governing relative chain, there exist
\[
       z_{\mathcal E}\in D_n\otimes A_n
       \quad\text{and}\quad e_{\mathcal E}\in F
\]
such that
\begin{align}
       \partial_n z_{\mathcal E}&=B_{\mathcal E},
                                                        \tag{E2}\label{eq:E2}\\
       \src_n(z_{\mathcal E})&=P_{\mathcal E}+e_{\mathcal E},
                                                        \tag{E1}\label{eq:E1}\\
       \ev(e_{\mathcal E})&\in\gamma_{n+1}(L),
       \qquad
       \operatorname{coord}_C\bigl(\ev(e_{\mathcal E})\bigr)=0,
                                                        \tag{E0}\label{eq:E0}
\end{align}
where $C=\gamma_{n+1}(L)$ is the cyclic terminal layer and
$\operatorname{coord}_C$ is its chosen cyclic coordinate.
\end{lemma}

This is exactly the formulation consumed by the current Step--7 assembly.
It may equivalently be stated with a genuine relation in place of the error:
there are $z_{\mathcal E}$ and $r_{\mathcal E}\in\Rel$ satisfying
\eqref{eq:E2} and
\[
       \src_n(z_{\mathcal E})=P_{\mathcal E}+r_{\mathcal E}.
\]
Indeed, the coordinate map on $C$ is injective, so \eqref{eq:E0} implies
$\ev(e_{\mathcal E})=0$, hence $e_{\mathcal E}\in\Rel$.  Conversely, a
genuine relation gives \eqref{eq:E0} with zero evaluation.  No stronger
claim about a separated homogeneous component of a relation is involved.

\section{The exact point at which the present verification stops}

All the following ingredients have already been verified independently.
\begin{enumerate}
\item Proper-subset collection leaves exactly the signed edge
      \(B_{\mathcal E}\).
\item Collecting each literal word \(\iota(\rho_e)X_e\) while retaining
      its full relation gives a genuine terminal chain, but its boundary is
      \(B_{\mathcal E}\) minus a terminal component read.
\item The unbracketed factor-two and factor-one reads of a triangular
      weight-one relation depend only on its complete grouped Smith head.
\item All nonexceptional rows and every correction whose full relation
      begins in weight at least two have genuine terminal chains with exact
      boundary and primitive formulas.
\item The horizontal and vertical reads of every intermediate diagonal
      agree by
      \[
          \Phi_k(f_k\chi_k)=T_k(\partial\chi_k),
      \]
      and the governing PBW coefficients make the right-hand side zero.
\item Point~6 annihilates the terminal coordinate of every genuine terminal
      source cycle.
\end{enumerate}

The missing implication is the passage from items 2--5 to one pair
\((z_{\mathcal E},e_{\mathcal E})\) satisfying
\eqref{eq:E2}--\eqref{eq:E0}.  In other words, the unverified object is not
merely a chain with the right boundary.  It is a chain with the right
boundary whose \emph{complete source primitive} has the right value modulo a
whole relation, or equivalently modulo a terminal-layer error whose cyclic
coordinate has been proved to vanish.

\section{Why the two tempting local arguments do not prove the lemma}

First, subtracting the already known proper-subset chain from the already
known total exceptional boundary chain produces \eqref{eq:E2} formally, but
it is circular.  When that chain is inserted into the final assembly, it
cancels the very complete chain whose closedness is being established.  It
does not provide an independent terminal-source occurrence trace.

Second, write a triangular relation as
\(\rho=\rho_1+\rho_{>1}\).  In a context of weight \(n\), the higher tail is
above the nilpotency cutoff and the replacement of \(\rho\) by \(\rho_1\) is
valid.  In a context of weight \(n-1\), only the prefix equality
\[
 \pi_n(c(\rho))=\pi_n(c(\rho_1))
\]
is valid.  In still lower contexts the omitted terms are nonzero.  Those
terms are exactly the intermediate horizontal and vertical diagonal edges.
Consequently a cellwise Smith-head replacement cannot prove
\eqref{eq:E1}.

\section{How the original proof proves Lemma~\ref{lem:missing}}

The required argument is the global marked-row Stokes calculation, restricted
only after its complete external boundary has been formed.

Start with the \emph{entire} occurrence-labelled initial family $I$ of the
governing relative chain, and mark within its trace the occurrences whose
upper external edge belongs to $\mathcal E$.  Do not first discard the other
initial rows: some descendants of an exceptional row meet correction rows
coming from lower contexts, and those rows supply its diagonal partners.
Retain every full relation as a whole.  Perform ordinary adjacent
interchanges, full-relation interchanges, and successive weight truncations.
A relation interchange creates the full relation $[x,\rho]\in\Rel$; it never
creates a bare homogeneous pseudo-relation.  Every occurrence created by one
replacement is kept until it is either used as the input of a later
replacement or remains on the terminal frontier.  Only after all internal
and diagonal incidences have been paired may the already certified
nonexceptional external families be removed from the equality.

The occurrence ledger has the exact telescoping identity
\[
 I-F_N=\sum_s\bigl(\text{input}(s)-\text{signed outputs}(s)\bigr).
\]
Reading the terminal factor-two rows that remain after this global
cancellation gives a genuine element of $D_n\otimes A_n$.  Its uncancelled
upper edge is precisely $B_{\mathcal E}$, giving \eqref{eq:E2}.  For the
one-factor read, all internal occurrences cancel before any top projection is
applied.  The part of the remaining external boundary relevant to the
exceptional edge is:
\begin{center}
\begin{tabular}{c|c}
terminal family & signed one-factor contribution\\ \hline
full one-factor relations & an element of \(\Rel\)\\
horizontal diagonal at \(k\) & \(-\Phi_k(f_k\chi_k)\)\\
vertical diagonal at \(k\) & \(+T_k(\partial\chi_k)\)\\
terminal factor-two rows & \(\src_n(z_{\mathcal E})\)\\
initial exceptional edge & \(-P_{\mathcal E}\)
\end{tabular}
\end{center}
with $2\leq k<n$.  The horizontal and vertical terms cancel by the
already established identity
$\Phi_k(f_k\chi_k)=T_k(\partial\chi_k)$, with opposite boundary
orientations.  The governing coefficient theorem also makes their total
common value zero.  After the surviving full relations have been grouped,
the unprojected remainder is the element $e_{\mathcal E}$ in
\eqref{eq:E1}.  The external-edge classification places its evaluation in
$C$, and the evaluated closed-square identity says exactly that its
$C$-coordinate is zero.  This proves \eqref{eq:E0}; injectivity of the cyclic
coordinate then permits it to be replaced by a genuine relation.

The crucial order is: sum the complete occurrence ledger, cancel its
internal incidences, pair every lower-context diagonal, remove the already
certified external families, group the surviving full relations, and only
then take the terminal coordinate.  The phrase ``the exceptional trace''
must refer to this resulting signed subledger, not to an independently
collected family obtained by deleting the other rows at the start.

\section{Why this single lemma is sufficient}

Add \(z_{\mathcal E}\) to the already constructed proper-subset tail chain.
Equation \eqref{eq:E2} gives the complete exceptional component boundary,
and \eqref{eq:E1}--\eqref{eq:E0} give its primitive modulo a zero-coordinate
terminal error (equivalently a genuine relation).  Adding the existing
nonexceptional and derived-tail corrections then gives a chain whose boundary
is the complete raw cutoff factor-two symbol and whose primitive is the
complete raw cutoff primitive modulo a zero-coordinate terminal error.
Adding the already constructed opposite complete factor-two chain produces a
genuine terminal source cycle.

The primitive of that cycle has the governing top coordinate.  Point~6 says
that the coordinate of a genuine terminal source cycle is zero.  Hence the
governing element is zero.  This is the unconditional conclusion of Step~7;
the reduced theorem and the final reduction in Step~8 then require only the
Reduction Property and the Passi--Sicking Property.

\end{document}
